% !TeX root = ../papers/nc-p3.tex
\section{非正規三次の導手繊維と固定した残余母線}
\label{sec:cofactor-cubic-conductor-base}
\begin{lemma}\label{lem:cofactor-sectional-adjoint}
定理\ref{thm:cofactor-sectional-graph-genus}の一般切断グラフ $G$ はGorensteinで、
$\omega_G=\mathcal O_G(2L-H)$ である。$L,H$ はそれぞれ核と入力の超平面である。
\end{lemma}
\begin{proof}
平面表示 $\mathcal O(-2)^4\to\mathcal O^4\to M_H$ の対称性から平面八次 $D_H$ 上で
$\mathcal Hom(M_H,\omega_{D_H})=M_H(-1)$。有限双有理押出し $M_H=h_*\mathcal O_G(L)$ の
有限双対は $h_*\omega_G(-L)$ である。同型は一般点で分数体の乗法、
従って両端の $\mathcal O_G$ 加群についても線形である。有限射で戻すと表示を得る。
\end{proof}

\begin{theorem}\label{thm:cofactor-cubic-conductor-base}
$d=18$ で非正規三次核曲面が錐でなければ矛盾する。
特異平面三次を底とする三次錐でも、母線の入力次数二の場合は生じない。
\end{theorem}
\begin{proof}[基点を数える]
算術種数三の $J=G$ で逆写像が二重直線 $T$ の一般点に降りない場合だけを扱う。
非錐なら定理\ref{thm:cofactor-nineteen-surface}の正規化 $S_0=\mathbb F_1$, $L=E+2F$ を使う。
$A=mE+(6-m)F$ が逆系で、$A^2=12m-3m^2\ge8$, $AE\ge0$ より $m=1,2,3$。
$T$ の逆像 $Z_0$ は $E+F$ のCartier因子で、$\rho:Z_0\to T$ は平坦有限二重被覆
（既約切断または切断と一本の母線）である。
$K$ の関数とは、正規化の関数で $Z_0$ への制限が $T$ の関数になるものである。
実際 $T$ のイデアルは負でない母線座標を掛けた $r$ の二生成表示で $\mathcal O(-Z_0)$ に引き戻される。
上の部分代数には $\mathcal O_K$ が入り、両者のHilbert多項式は等しい。
正規化と三次曲面のEuler数差が $t$、$Z_0$ と $T$ のEuler数差も $t$ だからである。

錐では平面三次の正規化で $S_0=\mathbb F_3$, $L=E+3F$、逆系は $mE+6F$, $m=1,2$ である。
ここでは $m=2$ を扱うので $A=2L$、一般の固有切断 $J_0$ は $E$ を避ける。
頂点外の導手の表示は同様である。底の特異平面三次は正規化余長一で、導手の逆像は
相異なる二点または長さ二の一点、そこへの制限が定数である関数が元の環となる
（二枝の節点または一枝の尖点）。底のアフィン開集合上で錐母線方向の座標を加えると
$Z_0$ は二本または二重の母線で、頂点外の $T$ に平坦二対一である。

一般逆切断 $J_0=(c=0)$ の像 $J$ への正規化方向の余長を計算する。
$T$ の局所環で二次代数 $\rho_*\mathcal O_{Z_0}$ の基底を $1,w$ とし、束を自明化して
$c|_{Z_0}=c_0+c_1w$ と書く。$J_0$ と $J$ の構造層の差は
$\mathcal O_{Z_0}/(\mathcal O_T+c\mathcal O_{Z_0})$ で、長さは $(c_0,c_1)$ の共通零次数である
（$w^2=aw+b$ として $w$ の係数だけを読む）。
逆写像が降りないとは四切断の $Z_0$ 制限が $T$ の一般点で階数二を持つことに等しい。
実際一般点で単元の一切断で割ればこの階数が一のとき全比が元の局所環に入る。
従って一般 $c$ の共通零因子は固定因子 $\Delta$ であり、その引戻し $P=Z_0\times_T\Delta$ は
全入力切断の基点スキームに入る。固定共通零を除いた後、一般で階数二の条件は
底の各一般点で余次元二、残る有限個の点でも非零切断を選べるためである。
$\deg\Delta=p_a(J)-p_a(J_0)=3-p_a(J_0)$。非錐の $m=1,2,3$ では
$(p_a(J_0),AZ_0)=(0,5),(2,4),(1,3)$ である。
$2\deg\Delta\le AZ_0$ により第一、第三の場合は消える。

残る両モデルで $A=2L$, $K_{S_0}+A=F$。$P$ は一点 $q\in T$ 上の長さ二繊維である。
$J_0\to J$ の余長一は全て $q$ にあり、導手は $J$ の極大イデアルで
$J_0$ では $\mathcal O_{J_0}(-P)$ となる。実際 $q$ のスキーム逆像は $P$、
両曲線はGorensteinなので有限双対で導手は $J_0$ 上の可逆イデアルとなる。
$G=J$ の入力比が $J_0$ に引き戻せるので、逆基点 $Z\subset S_0$ は二生成の完全交叉で長さ四
（一般切断を法として主イデアルとなる）。一般 $J_0$ 上で $Z$ はCartierかつ
$H|_{J_0}=A|_{J_0}(-Z)$ である。前補題と
$\omega_{J_0}=\nu^*\omega_J(-P)$ を使えば
\[
        \mathcal O_{J_0}(Z-P)=\omega_{J_0}=\mathcal O_{J_0}(F).
\]
$Z-P$ は面上でも固定した長さ二スキーム $Z_1$ であり、イデアル商 $I_Z:I_P$ で与えられる。
種数二の標準切断は母線の二切断で尽くされるので、$Z_1$ は一本の固定母線に含まれる。
\end{proof}

\begin{proof}[十六個を超える三点繊維]
$Y=\operatorname{Bl}_Z S_0$ とし、入力射を $f$、$N'=f_*\mathcal O_Y(L)$ とする。
表示の余核を $M$ とすれば、有限入力開集合で $N'/M$ は零か純一次元Cohen--Macaulayである
（局所完全交叉面の有限直像と $M$ の深さを使う）。
一般入力平面による切断では $Y$ 上の切断が $J_0$ と同型となる。
実際二生成イデアルのブローアップのチャートで一般切断を零とするとその比の座標が一意に決まる。
従って $N'/M$ の一般平面切断は $q$ からの一点で長さ一、支持は入力直線 $\Lambda$ である。
有限な開集合ではこの商は滑らかな直線上の線束となる。実際重複度一かつ純なので
直線イデアルの積は有限支持で零、直線上では階数一torsion-freeである。
$P$ の各台点上の例外直線は入力でこの同じ $\Lambda$ へ同型に写る。
一般入力平面でそれらが全て $q$ の上の一点の入力像を持つことと、相対超平面次数一による。

面上の重複度平方和は $A^2-8=4$。一二重基点なら $P,Z_1$ ともその点に支持される。
入力が潰す全曲線の像は $\Lambda$ にある。$\mathbb F_1$ で $F,E$ 以外の整曲線
$aE+bF$ は $b\ge a>0$、$A$ 次数 $2(a+b)$ に対し重複度は $a$ 以下なので縮約されない。
$\mathbb F_3$ でも $b\ge3a$、$A$ 次数 $2b$ で同様。
縮約される母線は基点を通り、$\mathbb F_1$ の負切断も縮約されるなら基点を通る。
$\mathbb F_3$ の負切断は基点を通らないが、基点を通る母線の固有変換と交わり、
その母線の入力次数は $2-2=0$ である。従って $\Lambda$ 外では全繊維が有限で $M=N'$。
直線上の印は高々三個なので少なくとも十七個の孤立三点繊維を得る。

四単純基点の場合は近接を含め各台点で自由な鎖であり、長さ $r$ の鎖のブローアップは
$(v,u^r)$ のブローアップである。それぞれの例外直線上で $Y$ が特異になり得る点は高々一つである
（局所式 $vX=u^r Y'$）。入力の正次元繊維は高々三入力点である。
$\mathbb F_3$ では負切断と二基点を要する母線高々二本だけである。
$\mathbb F_1$ では負切断と母線以外の収縮には $2(a+b)\le4a$ が必要なので $b=a$、
全基点での重複度が $a$ となる。平面に負切断を潰すと次数 $a$ の曲線であり
重複度 $a$ の点を持つなら直線である。従って $E+F$ の一員が四基点全てを通る場合だけで、
このとき負切断は交わらず、各母線との交点も一以下なので他の縮約はない。

$P$ が二点なら、それらは相異なる母線にある（各母線で核射は同型）。
$Z_1$ は一本の母線上なので二点の少なくとも一方の鎖の長さは一である。
従って $\Lambda$ 上に、$Y$ が全点で滑らかで $\Lambda$ に同型に写る一例外直線がある。
$P$ が一点の場合もその例外直線を使い、$Y$ がそこで特異な高々一入力点を除く。
この場合は正次元繊維の入力点も高々二つで済む。長さ二の $P$ の接方向は核への微分で消え、
母線方向は消えないので異なる。$Z$ の台はこの一点と $Z_1$ の固定母線上に限られ、
その母線以外で二基点を含む母線はない。負切断以外の水平縮約は上記の全四点を通る
$E+F$ の場合だけである。

以上の除外点の外で $\Lambda$ 上の各印も少なくとも一つの孤立三点繊維を与える。
その局所代数を詳述する。$x$ をこのような印とし、$n=\dim N'/\mathfrak m_x N'$ を有限繊維長とする。
$N'/M$ は直線束ゆえ $n\le\dim(M/\mathfrak m_x M)+1=4$。
直線上の一般点で $n=2$ である。二台の $P$ なら二例外点で一般平面の $J_0$ は滑らかであり、
一台でも単純鎖の一般 $J_0$ は滑らかである。$J_0\to J$ で導手イデアルは主で商長二、
入力はそこで階数三なので正規化の入力繊維長は二となる。
商への値に寄与する局所成分は $P$ 上の例外点の成分だけである。
他の成分から直線束への射は $\Lambda$ の一般点で零ゆえ零である
（完備化またはhensel化し、有限射の各台の直和分解を用いる）。

$n\le2$ を排除する。二台なら $n=2$ であり両台の長さは一。
商の一般切断は二点の値の差であるから、有限開集合で $M$ は
二例外直線上の値の差の核に等しい。核座標の値が同じなので包含があり、差の評価は全射、
商がともに直線束で一般同型なので深さ二の延長で等号である。
各例外直線上の $L$ は固定核点の束で、値の同一視は全点で可逆である。
滑らかな方の台では長さ一なので入力への射は局所的に滑らかな面の閉埋込みとなる。
その面の局所方程式は直線を含み一次項非零である。従って
$\operatorname{Tor}_1(N',k_x)\to\operatorname{Tor}_1(\mathcal O_\Lambda,k_x)$ は階数一以上である
（その方程式の直線イデアルの二生成元に関する一次係数を取る）。
完全列から $\dim M/\mathfrak m_xM\le n-1+2-1=2$ で矛盾する。
一台の場合は上の商に寄与する台で $Y$ は滑らかで、一般の $x$ を通る平面は例外直線を横断し
そこで滑らかな曲線を切る。この平面で商は長さ一、包含は単射のままで、当該成分で
$M$ の切断加群は滑らかな曲線環の有限余長部分加群となる。
その最小生成数は入力平面の極大イデアルによる重複度以下である。
この重複度は滑らかな曲線環と同じで、その環では極大イデアルの引戻しが主なので
ちょうど当該有限繊維長に等しい。他の局所成分では $M=N'$。
従ってこの場合も $3\le n$ が必要となる。

$n=3$ なら繊維自身、$n=4$ なら上記の滑らかな例外点上の局所繊維長を一減らして
他の台の繊維を全て取る。その点の繊維は一般平面切断の滑らかな曲線に含まれて曲線的であり、
この配分で長さ三の部分スキームは唯一である。
得られた点は孤立する。有限開集合で $\Lambda$ の外では三点以上の繊維の台が有限、
$\Lambda$ 自身の一般繊維長は二であり、閉繊維内でも長さ配分は開閉だからである。
これで各場合に二十印のうち高々三入力点を除いて孤立三点繊維を得た。

一方、局所完全交叉長さ四の基点と入力切断は平滑化できる。
一般切断 $J_0$ の算術種数二と入力次数八より $H^1(I_Z A)=0$。
$\mathbb F_3$ は $\mathcal O\oplus\mathcal O(3)$ を
$\mathcal O(1)\oplus\mathcal O(2)$ へ動かして $(\mathbb F_1,L)$ に移せる。
一般には四相異点の $(\mathbb F_1,2L)$ のブローアップである。
各四母線に一点ずつ置けば母線直像の三直和次数 $4,3,2$ から四条件を引く。
二つの直和端の座標点で $4$ の因子に三条件、$2$ の因子に一条件を課せば残りは $1,3,1$。
従って一般には正次数スクロールの $2\xi-2F$ 因子への埋込みであり、
一般の四切断の三点繊維は有限である。平面型 $4h-2e_0-\sum_{i=1}^4e_i$ の同伴束積分は十六である。
十七以上の孤立点に反する。\path{computations/verify_cofactor_cubic_conductor.py}
は導手長とこの積分を検算する。
\end{proof}
